% Einheit 4 von 4 – Ausklammern (bank/terme/e4.jsonl, zusammenbau v0.1)
%% TODO Zweigzeile Teil 1: was hier gelernt wird (Satz in Schülersprache aus dem Einheitstitel) – Einheit 4
\einheitenkopf[e4]{Einheit 4 von 4 · Ausklammern}
\zweigzeile{ab Kl. 8 (am Gymnasium ab Kl. 7) · keine P10-Aufgabe}
\setcounter{aufgabe}{19}

%% TODO Titel: Ich-kann-Satz fehlt in der Bank (Platzhalter: Kettenname) – Nr. 20
%% TODO Anweisung: ein Satz über der ersten Teilaufgabe fehlt in der Bank – Nr. 20
\begin{aufgabe}{Ausklammern}
%% TODO \rechenplatz{n}: Schrittzahl des Grundfalls fehlt in der Bank (nur bei fester Schreibform, 2.3 b)
\begin{teile}
\teil Klammere den größten gemeinsamen Faktor aus: $4x + 8$ \leerfeld
\teil Klammere den größten gemeinsamen Faktor aus: $x^2 + 5x$ \leerfeld
\teil Klammere den größten gemeinsamen Faktor aus: $6x^2 + 15x$ \leerfeld
\teil Klammere den größten gemeinsamen Faktor aus: $6x + 6$ \leerfeld
\teil Klammere den größten gemeinsamen Faktor aus: $3x^2 + 6x + 15$ \leerfeld
\teil Klammere aus und mache die Probe: $6x + 30$
\teil Finde den Fehler und rechne richtig: \rechnung{5x + 20 &= 5 \cdot (x + 20)}
\teil Klammere so weit wie möglich aus und mache die Probe: $3a^2 + 6ab + 9a$
\end{teile}
\end{aufgabe}

%% TODO Titel: Ich-kann-Satz fehlt in der Bank (Platzhalter: Kettenname) – Nr. 21
%% TODO Anweisung: ein Satz über der ersten Teilaufgabe fehlt in der Bank – Nr. 21
\begin{aufgabe}{Ausklammern – Fehler finden}
\begin{teile}
\teil Finde den Fehler und rechne richtig: \rechnung{12x - 8 &= 4 \cdot (3x + 2)}
\end{teile}
\end{aufgabe}
